\chapter{کار با فایل و سیستم}
\label{ch:files}

تا اینجا داده‌ها را درونِ برنامه می‌ساختیم و خروجی را روی صفحه چاپ می‌کردیم.
اما برنامه‌های واقعی با دنیای بیرون هم در ارتباط‌اند: فایل می‌خوانند و
می‌نویسند، آرگومان‌های خطِ فرمان می‌گیرند و با سیستم‌عامل کار می‌کنند. در این
فصل این پل‌ها را می‌سازیم.

\section{نوشتن و خواندنِ فایل}

ماژول‌های \code{io} و \code{os} ابزارهای کار با فایل را فراهم می‌کنند. این مثالِ
واقعی یک فایل می‌نویسد، آن را بازمی‌خواند و سپس پاک می‌کند:

\begin{salamen}
package main
import (
    "os"
    "io"
    "str"
)
func main:
    path str = "io_example_out.txt"
    io.WriteFile(path, "the quick brown fox\njumps over\nthe lazy dog\n")
    text str = io.ReadFile(path)
    lines Vector<str> = str.Split(str.Trim(text), "\n")
    io.Println("read lines:")
    __ str.Join(lines, " | ")
    os.Remove(path)
    lines.free()
end
\end{salamen}

\begin{note}
این نمونه به انگلیسی نوشته شده تا با کدِ واقعیِ کتابخانه هماهنگ باشد. به یاد
داشته باشید که \code{\_\_} (دو زیرخط) شکلِ کوتاهِ \code{چاپ}/\code{println}
است؛ پس \code{\_\_ str.Join(...)} یعنی «نتیجه را چاپ کن». همین برنامه را
می‌توان با واژگانِ فارسی نیز نوشت.
\end{note}

سه عملِ کلیدیِ فایل را اینجا می‌بینیم:
\begin{itemize}
  \item \code{io.WriteFile(مسیر, محتوا)}: یک رشته را در فایل می‌نویسد (و اگر
        فایل نبود، می‌سازد).
  \item \code{io.ReadFile(مسیر)}: همهٔ محتوای فایل را به‌صورتِ یک رشته
        برمی‌گرداند.
  \item \code{os.Remove(مسیر)}: فایل را پاک می‌کند.
\end{itemize}

پس از خواندن، با \code{str.Split(str.Trim(text), "\textbackslash n")} متن را به
خط‌هایش شکستیم همان \code{بشکاف}/\code{Split} که در فصلِ ۱۵ دیدیم. این الگو
--- بخوان، با \code{Split} بشکن، پردازش کن ستونِ فقراتِ کار با فایل‌های
متنی است.

\begin{warn}
خواندن و نوشتنِ فایل می‌تواند شکست بخورد: فایل ممکن است وجود نداشته باشد، یا
اجازهٔ دسترسی نباشد. در برنامه‌های جدی، پیش از خواندن با \code{os.Exists(مسیر)}
بررسی کنید که فایل هست، و همان‌طور که در فصلِ ۱۳ گفتیم، حالتِ خطا را آگاهانه
مدیریت کنید.
\end{warn}

\section{آرگومان‌های خطِ فرمان}

وقتی برنامه‌ای را اجرا می‌کنید، می‌توانید همان لحظه به آن داده بدهید. مثلاً:

\begin{salamen}[language={}]
./app.exe سلام دنیا
\end{salamen}

اینجا «سلام» و «دنیا» دو \emph{آرگومان}اند. با \code{os.Args()} می‌توانید
آن‌ها را در برنامه بگیرید که یک بردار از رشته‌ها برمی‌گرداند:

\begin{salamen}
package main
import ("os" "io" "str")

func main:
    a Vector<str> = os.Args()
    io.Print("called with ")
    io.Print(str.FromInt((a.len() - 1) as i64))
    io.Println(" argument(s)")
    mut i = 1
    while i < a.len():
        io.Println(a.get(i)[0])
        i = i + 1
    end
    a.free()
end
\end{salamen}

\begin{note}
نخستین عنصرِ \code{os.Args()} (یعنی اندیسِ ۰) همیشه \emph{نامِ خودِ برنامه}
است؛ به همین دلیل حلقه را از اندیسِ ۱ آغاز کردیم و برای شمارشِ آرگومان‌های
واقعی، یکی از طول کم کردیم (\code{a.len() - 1}). این قراردادی است که در بیشترِ
زبان‌ها مشترک است.
\end{note}

\section{متغیرهای محیطی}

سیستم‌عامل اطلاعاتی را به‌صورتِ \emph{متغیرهای محیطی} نگه می‌دارد مانندِ
نامِ کاربر یا مسیرِ خانه. برنامه‌ها می‌توانند این‌ها را بخوانند. ماژولِ
\code{os} تابعِ \code{Env} را برای این کار دارد:

\begin{salamen}
package main
import ("os" "io")

func main:
    user str = os.Env("USERNAME")
    io.Println(user)
end
\end{salamen}

\begin{tip}
متغیرهای محیطی جای خوبی برای نگه‌داشتنِ تنظیماتِ حساس (مانندِ کلیدهای دسترسی)
هستند، چون درونِ کدِ شما نوشته نمی‌شوند و در هر سیستم می‌توانند متفاوت باشند.
این یک عادتِ مهمِ امنیتی است: \emph{رمزها و کلیدها را هرگز مستقیم در کد ننویسید.}
\end{tip}

\section{تشخیصِ سیستم‌عامل در زمانِ کامپایل}

گاهی برنامه باید روی ویندوز، مک و لینوکس کمی متفاوت رفتار کند مثلاً
جداکنندهٔ مسیرها در ویندوز \code{\textbackslash} و در لینوکس \code{/} است. سلام
سینتاکسِ جداگانه‌ای برای این کار ندارد همان دستورِ \code{اگر}/\code{وگرنه}/
\code{پایان} را که می‌شناسید، در سطحِ بالای فایل (بیرونِ هر تابعی) به کار
می‌برید، و کامپایلر آن را در \emph{زمانِ کامپایل} حل می‌کند. این مثالِ واقعی را
ببینید:

\begin{salamen}
if SALAM_OS_WINDOWS:
    const PATH_SEP = "\\"
    const SO_EXT   = ".dll"
else if SALAM_OS_MAC:
    const PATH_SEP = "/"
    const SO_EXT   = ".dylib"
else if SALAM_OS_UNIX:
    const PATH_SEP = "/"
    const SO_EXT   = ".so"
end

func main:
    if SALAM_OS_WINDOWS:
        __ "Running on Windows"
    else:
        __ "Running on a Unix-like system"
    end
    __ "Path separator:", PATH_SEP
end
\end{salamen}

پرچم‌هایی مانندِ \code{SALAM\_OS\_WINDOWS}، \code{SALAM\_OS\_MAC} و
\code{SALAM\_OS\_UNIX} را خودِ کامپایلر بر اساسِ سیستمِ مقصد تعیین می‌کند و
همیشه یک مقدارِ \code{true}/\code{false} مشخص دارند، پس شرطِ بالا همیشه قابلِ
حلِ در زمانِ کامپایل است.

\begin{deep}[نگاهی به درون: چه زمانی \code{اگر} در زمانِ کامپایل حل می‌شود؟]
یک \code{اگر} در \emph{سطحِ بالای فایل} (بیرونِ هر تابعی) همیشه باید در زمانِ
کامپایل قابلِ حل باشد چون در آن‌جا هیچ متغیرِ زمانِ اجرایی وجود ندارد؛
شاخهٔ ناخواسته اصلاً واردِ برنامهٔ نهایی نمی‌شود و حتی بررسیِ نوع هم روی آن
اجرا نمی‌شود. به همین دلیل برای کدی که فقط روی یک سیستم \emph{معنا} دارد
(مثلِ فراخوانیِ یک کتابخانهٔ مخصوصِ ویندوز) ضروری است: نمی‌خواهیم آن کد روی
لینوکس حتی کامپایل شود.

یک \code{اگر} \emph{داخلِ یک تابع} برعکس است: معمولاً در زمانِ اجرا تصمیم
می‌گیرد (هر دو شاخه در برنامهٔ نهایی حضور دارند)، مگر این‌که شرطش دقیقاً به
یکی از همین پرچم‌های کامپایلری مانندِ \code{SALAM\_OS\_WINDOWS} اشاره کند
در آن صورت کامپایلر همان‌جا هم آن را در زمانِ کامپایل حل می‌کند. یک شرطِ
معمولی مانندِ \code{اگر سن > ۱۸} هیچ ربطی به این مکانیزم ندارد و همیشه در
زمانِ اجرا بررسی می‌شود.
\end{deep}

می‌توانید پرچم‌های دلخواهِ خود را هم هنگامِ ساخت تعریف کنید:

\begin{salamen}[language={}]
salam build app.salam -DDEBUG --output=app.exe
\end{salamen}

و در کد با \code{اگر defined(DEBUG): ... پایان} به آن پاسخ دهید روشی عالی
برای گذاشتنِ پیام‌های اشکال‌زدایی که فقط در نسخهٔ توسعه فعال‌اند. برخلافِ
پرچم‌های سیستم‌عامل که همیشه تعریف‌شده‌اند، پرچم‌های دلخواه ممکن است اصلاً
پاس داده نشده باشند؛ \code{defined(...)} دقیقاً همین حالت را بی‌خطا بررسی
می‌کند.

\section{خلاصهٔ فصل}

\begin{itemize}
  \item با \code{io.WriteFile} و \code{io.ReadFile} فایل می‌نویسیم و می‌خوانیم.
  \item \code{os.Remove} و \code{os.Exists} برای مدیریتِ فایل‌اند.
  \item \code{os.Args()} آرگومان‌های خطِ فرمان را می‌دهد (اندیسِ ۰ نامِ
        برنامه است).
  \item \code{os.Env(...)} متغیرهای محیطی را می‌خواند.
  \item همان \code{اگر}/\code{وگرنه}/\code{پایان}، وقتی در سطحِ بالای فایل
        به کار رود، کد را بر اساسِ سیستم‌عامل در زمانِ کامپایل انتخاب می‌کند.
\end{itemize}

\section{تمرین‌ها}

\exercise برنامه‌ای بنویسید که یک پیام را در فایلی بنویسد و سپس آن را بخواند و
چاپ کند.

\exercise برنامه‌ای بنویسید که نام شما را به‌عنوانِ آرگومان بگیرد و «سلام،
<نام>» چاپ کند.

\exercise فایلی متنی بسازید که چند خط دارد، آن را بخوانید و تعدادِ خط‌ها را
بشمارید.

\exercise با کمکِ \code{@if}، برنامه‌ای بنویسید که روی هر سیستم‌عامل پیامِ
متفاوتی چاپ کند.

\exercise برنامه‌ای بنویسید که چند عدد را به‌عنوانِ آرگومان بگیرد و مجموعشان را
حساب کند (راهنمایی: از \code{به\_صحیح}/\code{ToInt} استفاده کنید).
